% Owner: E. M4 / H5, input from results_evolution.tex.
% Values: results/e_summary.csv and results/e/m4_*; figures: scripts/plot_e.py.
\begin{figure*}[t]
\centering
\includegraphics[width=\linewidth]{e_m4_ps_timeseries.pdf}
\caption{Snapshot communities in PS (15-min windows, $\gamma=1$). Rows: NMI of each
snapshot partition with the classes, with the previous window and with the static
partition. Blue: real; orange: P1 (mean $\pm$ sd, 20 samples); green dotted: same window
partitioned with another seed (noise floor).}
\label{fig:m4-ps}
\end{figure*}

\begin{figure}[t]
\centering
\includegraphics[width=\columnwidth]{e_m4_cm_timeseries.pdf}
\caption{Snapshot communities in CM (7-day windows, $\gamma=1$); encoding as in
\cref{fig:m4-ps}. CM has no ground-truth labels.}
\label{fig:m4-cm}
\end{figure}

\begin{figure*}[t]
\centering
\includegraphics[width=\linewidth]{e_m4_sensitivity.pdf}
\caption{Sensitivity of H5 to window width and resolution $\gamma$. Top two rows: real
minus P1 mean of the window-averaged NMI with the static partition and with the previous
window (negative: real snapshots agree less). Bottom: communities per real snapshot.}
\label{fig:m4-sens}
\end{figure*}

\paragraph{Dynamic vs.\ static communities.}
In PS (\cref{fig:m4-ps}) the snapshot communities mostly match the classes outside the
midday break (NMI $\approx 0.85$--$0.95$), as does the static partition (six communities,
NMI 0.86 with the classes). Between about 12:00 and 14:00 on both days the NMI with the
classes falls to $\approx 0.5$, and the NMI with the static partition falls just as far.
Under P1 both curves stay flat (class NMI $0.86\pm0.002$), so the drop comes from when
contacts happen, not from who is in contact. It is not algorithmic noise either: two
Leiden runs on the same window agree with NMI 0.98 on average. Averaged over the day,
real snapshots agree with the static partition less than P1 snapshots (0.72 vs.\
$0.94\pm0.003$).

CM behaves differently (\cref{fig:m4-cm}). Weekly snapshots agree only weakly with each
other and with the static partition (NMI 0.52 and 0.34), and P1 gives almost the same
values (0.42 and 0.32). The snapshots are sparse (326 of 1{,}899 users active per week on
average), and during the busy first two months two runs of the same algorithm agree only
at NMI $\approx 0.65$--$0.7$. Where real and null differ, real snapshots are \emph{more}
similar to their predecessor than random ones, which points to persistent conversation
groups rather than reorganisation. The same contrast appears at the level of ties
(\cref{sec:r-activity}): the mean Jaccard overlap of consecutive edge sets is lower than
under P1 in PS (0.22 vs.\ 0.37, 1-h windows) but higher in CM (0.096 vs.\ 0.069, 7-day
windows).

One property of P1 matters for reading these numbers. Permuting timestamps spreads every
node's events over the whole observation period, so a P1 window contains almost all nodes
(240 vs.\ 175 active nodes per 15-min PS window) and resembles a sample of the static
graph. Part of the gap between real and P1 agreement with the static partition therefore
reflects \emph{who is active when}, not only \emph{who meets whom}. The lunch-time drop
in class NMI is not affected by this, since it compares real snapshots with external
labels.

\Cref{fig:m4-sens} repeats the comparison for three window widths and three resolutions.
In PS, real snapshots agree less with the static partition than P1 in all nine settings
(difference $-0.03$ to $-0.23$). In CM the difference stays within $\pm0.06$, and within
$\pm0.025$ except for 30\,days with $\gamma=0.5$. The comparison with the previous window
is not robust in PS: at 15 and 60\,min real snapshots change more than P1 ones, but at
5\,min the order flips, because a 5-min P1 snapshot is a small random sample of the static
graph and is itself noisy. Larger $\gamma$ gives more, smaller communities without changing
any of these signs.

\paragraph{Verdict on H5.}
\emph{Supported for face-to-face contacts, not supported for online messages.} In PS the
static partition recovers the classes but hides the daily mixing at lunch, which P1 does
not reproduce. In CM the static partition summarises weekly snapshots as well as it
summarises time-shuffled ones, and snapshots are, if anything, more persistent than random.

